\section{总结与延伸}

这节课最值得保留的不是模型清单，而是一个三层判断框架。第一层是 \textbf{representation/objective}：language token、continuous vision feature、discrete image code 与 noisy latent 各自保留什么信息。第二层是 \textbf{architecture/systems}：projection、experts、cross attention、ZeRO、parallelism 与 long context 怎样把表示映射到可训练硬件。第三层是 \textbf{data/evaluation}：instruction、preference、caption、GUI trace、video coverage 与 benchmark 如何定义模型最终学到的行为。

\subsection{三条主线的统一}

\begin{importantbox}{课程的最终统一图景}
\begin{enumerate}
  \item \textbf{LLM 提供通用序列底座：}scaling、alignment 与 systems 让大模型成为可复用接口。
  \item \textbf{Multimodality 暴露接口瓶颈：}视觉细节、高分辨率、二维空间与视频时间结构迫使模型改变 bridge、objective 和 compression。
  \item \textbf{Data 决定行为边界：}architecture 能提高可拟合范围，但 benchmark、annotation、recaption、verification 与 failure coverage 决定系统实际上学会什么。
\end{enumerate}
\end{importantbox}

\begin{warningbox}{最容易犯的三种归纳错误}
不要把同 pretraining loss 写成同能力定理；不要把一条 GUI trace 写成 agent reliability；不要把 diffusion 的 practical advantage 写成 autoregression 的不可能性证明。三种错误都来自把课堂证据层级抹平。
\end{warningbox}

\subsection{复现实验 checklist}

若要复现本讲任一模型，至少应同时发布 architecture version、data mixture/provenance、resolution/tokenization、training compute、inference steps、prompt/evaluation protocol 与 repeated-trial failures。只发布 headline score，会让 data、algorithm、architecture 三者无法区分。

\begin{lstlisting}[caption={跨 LLM/VLM/diffusion 的统一复现清单}]
pin(code_commit, model_checkpoint, tokenizer_or_vae)
document(data_sources, filters, synthetic_fraction, licenses)
report(train_flops, gpu_hours, parallelism, peak_memory)
report(input_resolution, context_length, sampling_steps)
freeze(prompts, decoding, benchmark_version, test_split)
publish(ablations, confidence_intervals, failures, costs)
\end{lstlisting}

\subsection{自测问题}

\begin{multicols}{2}
\footnotesize
\begin{enumerate}
  \setlength{\itemsep}{0pt}
  \setlength{\parskip}{0pt}
  \item GLM blank infilling 与纯 MLM、纯 autoregression 的条件信息有何不同？
  \item 为什么 scaling law 是预算分配器，而不是性能保证书？
  \item perplexity 为什么不能跨 tokenizer 直接比较？
  \item GQA 主要减少哪一类 inference memory？
  \item ZeRO-1/2/3 分别 sharding 什么？
  \item activation checkpointing 与 model checkpointing 有何区别？
  \item TP、PP、context parallelism 分别切哪一维？
  \item 为什么 long context 的主要交互成本可能在 prefill？
  \item DPO 没有独立 reward model，为何仍依赖 reward assumptions？
  \item CogQA 例子怎样体现 data/algorithm/architecture 可交换？
  \item BLIP-2 的 Q-Former 同时完成哪两类工作？
  \item LLaVA projection 简单，但为什么 OCR 仍可能差？
  \item CogVLM vision experts 想保护什么能力？
  \item CogAgent 为什么需要高分辨率旁路？
  \item 单条 web trace 不能证明哪些能力？
  \item Vary 与 GLM-4V 对 OCR bottleneck 的处理有何不同？
  \item image tokenizer 给 autoregressive generation 带来什么得失？
  \item diffusion 为什么通常比 batch-one AR decoding 更能利用 GPU？
  \item Relay Diffusion 为什么要改变跨 resolution 的 noise correlation？
  \item adaLN 与 cross attention 的 condition 路径有何不同？
  \item Sora-like recipe 中哪些是表示问题，哪些是 systems/data 问题？
  \item 为什么 annotated video pairs 不保证学到 physical rules？
\end{enumerate}
\end{multicols}

\subsection{拓展阅读}

\begin{multicols}{2}
\footnotesize
\begin{itemize}
  \setlength{\itemsep}{0.15em}
  \setlength{\parskip}{0pt}
  \item Du et al., \href{https://arxiv.org/abs/2103.10360}{GLM: General Language Model Pretraining with Autoregressive Blank Infilling}：理解目标函数统一。
  \item Rajbhandari et al., \href{https://arxiv.org/abs/1910.02054}{ZeRO}；Shoeybi et al., \href{https://arxiv.org/abs/1909.08053}{Megatron-LM}：理解状态分片与模型并行。
  \item Rafailov et al., \href{https://arxiv.org/abs/2305.18290}{Direct Preference Optimization}：理解 pairwise preference objective。
  \item Li et al., \href{https://arxiv.org/abs/2301.12597}{BLIP-2}；Liu et al., \href{https://arxiv.org/abs/2304.08485}{LLaVA}：比较 Q-Former 与 projection bridge。
  \item Wang et al., \href{https://arxiv.org/abs/2311.03079}{CogVLM}；Hong et al., \href{https://arxiv.org/abs/2312.08914}{CogAgent}：理解 vision experts 与 GUI high-resolution path。
  \item Ho et al., \href{https://arxiv.org/abs/2006.11239}{DDPM}；Peebles and Xie, \href{https://arxiv.org/abs/2212.09748}{DiT}；Esser et al., \href{https://arxiv.org/abs/2403.03206}{Stable Diffusion 3}：从 diffusion objective 到 Transformer denoiser。
  \item Snell et al., \href{https://arxiv.org/abs/2409.01236}{Same Pre-training Loss, Better Downstream}：作为课堂“loss 决定能力”强主张的 post-lecture 反例阅读。
\end{itemize}
\end{multicols}

\end{document}
